\family{Meliaceae}{Mahogany family}
\label{fam:meliaceae}

Trees or shrubs, the bark and wood often scented. Leaves alternate,
usually pinnate, without stipules. Flowers small, 4- or 5-merous, the
stamens joined into a tube. Fruit a capsule, a berry or a drupe; in
the timber genera treated here a woody capsule with winged seeds.
Generic limits follow \citet{pennington1975}. Three species in the
Arami Uplands, all valuable timber trees.

\section*{Key to the species}
\begin{florakey}
	\lead{1a}{Capsule cylindrical, pendulous, splitting from the base into
		5 valves; slash with a cedar-like scent}{\taxon{Entandrophragma
			cylindricum}}
	\lead{1b}{Capsule globose, erect, opening from the apex into 4 or 5
		valves}{2}
	\lead{2a}{Leaflets broadly obovate, up to 20~cm long; capsule
		6--8~cm across; forest tree}{\taxon{Khaya grandifoliola}}
	\lead{2b}{Leaflets oblong-elliptic, up to 10~cm long; capsule
		4--6~cm across; savanna tree}{\taxon{Khaya senegalensis}}
\end{florakey}

\species{Entandrophragma cylindricum}{(Sprague) Sprague}
\label{sp:MEL01}

\field{Synonyms} \sciname{Pseudocedrela cylindrica}{Sprague}

\field{Local and trade names} Trade: Sapele.

\field{Status} Native.

\field{Description} Deciduous tree to 50~m with short, blunt
buttresses. Bark grey-brown, flaking in small round scales; slash with
a cedar-like scent. Leaves paripinnate with 5--9 pairs of leaflets.
Flowers small, greenish-yellow. Capsule cylindrical, pendulous,
6--10~cm long, splitting from the base into 5 valves; seeds winged.

\field{Phenology} Flowering March to May; fruiting November to
January.

\field{Habitat} Lowland semi-deciduous forest.

\field{Uses} A reddish-brown timber with a striped figure, used for
veneer, furniture and joinery. Heavy logging has made large trees
scarce \citep{hawthorne2006,louppe2008}.

\field{Specimens examined} Olowo Stream (D3), 440~m, fr., 6~Jan 1936, \specimen{Harding}{1214}{OKF};
Igbo-Oba Forest Reserve (C3), 470~m, fr., 28~Jan 1937, \specimen{Harding}{1378}{OKF};
Olowo Stream (D3), 440~m, fr., 13~Nov 1938, \specimen{Harding}{1641}{OKF};
Ireke Ferry (D4), 360~m, fr., 27~Nov 2023, \specimen{Alade}{600}{AUH};
Ireke Ferry (D4), 360~m, fr., 8~Nov 2024, \specimen{Alade}{753}{AUH}.

\field{Distribution} Grid squares C3, D3, D4.

\begin{center}
	\distmap{C/3, D/3, D/4}
\end{center}

\species{Khaya grandifoliola}{C.DC.}
\label{sp:MEL02}

\field{Local and trade names} Yoruba: Oganwo. Trade: African
mahogany.

\field{Status} Native.

\field{Description} Deciduous tree to 40~m with buttresses. Bark
grey-brown, flaking in scales; slash red. Leaves paripinnate with
2--5 pairs of broadly obovate leaflets up to 20~cm long. Flowers
small, cream. Capsule erect, globose, 6--8~cm across, opening into 4
or 5 valves; seeds winged.

\field{Phenology} Flowering March to May; fruiting January to March.

\field{Habitat} Lowland semi-deciduous forest and forest--savanna
mosaic.

\field{Uses} A valuable mahogany timber for furniture and joinery. The
bitter bark is used in local medicine \citep{keay1989,louppe2008}.

\field{Specimens examined} Ireke Ferry (D4), 360~m, fr., 22~Mar 1935, \specimen{Harding}{1157}{OKF};
Okeri (C4), 420~m, fr., 11~Feb 1938, \specimen{Harding}{1534}{OKF};
Ewe Stream gallery forest (B3), 580~m, st., 19~Jul 1962, \specimen{Olaniyan}{356}{SVN};
Oke-Asa farmland (B4), 690~m, fr., 16~Jan 1967, \specimen{Olaniyan}{802}{SVN};
Olowo Stream (D3), 440~m, fr., 23~Mar 2020, \specimen{Alade}{252}{AUH}.

\field{Distribution} Grid squares B3, B4, C4, D3, D4.

\begin{center}
	\distmap{B/3, B/4, C/4, D/3, D/4}
\end{center}

\species{Khaya senegalensis}{(Desr.) A.Juss.}
\label{sp:MEL03}

\field{Synonyms} \sciname{Swietenia senegalensis}{Desr.}

\field{Local and trade names} Hausa: Madachi. Trade: Dry-zone
mahogany.

\field{Status} Native.

\field{Description} Tree to 30~m with a dense crown. Bark dark grey,
scaly. Leaves paripinnate with 3--6 pairs of oblong-elliptic leaflets
up to 10~cm long. Flowers small, cream. Capsule erect, globose,
4--6~cm across, opening into 4 valves; seeds winged.

\field{Phenology} Flowering March to May; fruiting January to
February.

\field{Habitat} Guinea savanna woodland, especially along
watercourses; also planted as a shade and avenue tree.

\field{Uses} Timber for furniture and building. The very bitter bark
is used in local medicine, and the branches are lopped for fodder in
the dry season \citep{arbonnier2004,orwa2009}.

\field{Specimens examined} Idofin (A3), 760~m, fr., 12~Jan 1962, \specimen{Olaniyan}{291}{SVN};
Asiri Crossing (A2), 740~m, fr., 27~Jan 1965, \specimen{Olaniyan}{525}{SVN};
Asiri Crossing (A2), 740~m, fr., 13~Jan 1967, \specimen{Olaniyan}{773}{SVN};
Aramiwo Escarpment, summit (B1), 900~m, fr., 27~Feb 2020, \specimen{Alade}{224}{AUH};
Asiri Crossing (A2), 740~m, fr., 11~Jan 2023, \specimen{Alade}{536}{AUH}.

\field{Distribution} Grid squares A2, A3, B1.

\begin{center}
	\distmap{A/2, A/3, B/1}
\end{center}